% !TEX TS-program = xelatex
% !TEX encoding = UTF-8 Unicode
% !Mode:: "TeX:UTF-8"
% Tailored for: AI应用工程师 (应届生/研发类) -- LLM application dev, fine-tuning/RAG-adjacent, effect evaluation

\documentclass{my-resume}
\usepackage{my-resume-zh}
\usepackage{linespacing_fix}
\usepackage{cite}

\begin{document}
\pagenumbering{gobble}

\photoheader{%
  \name{吴筱琦}
  \contactInfo{+86 13776020628}{yukiwxq@gmail.com}{现居城市：伦敦/苏州}{意向城市：上海/苏州/杭州 \textbar\ AI应用工程师}
}{../template/photo.jpg}{2.4cm}

\section{教育背景}
\entry{帝国理工学院 (QS 2)}{2025.09 - 2026.09}{计算生命科学研究型硕士}{伦敦}
\begin{itemize}[parsep=0.2ex]
  \item \textbf{课程内容}：R语言统计学、生态与进化数据科学、生物计算训练营等统计编程与数据科学核心课程，兼修 GIS、基因组学与生物信息学
\end{itemize}
\entry{帝国理工学院 (QS 2)}{2022.09 - 2025.07}{生物科学 本科}{伦敦}
\begin{itemize}[parsep=0.2ex]
  \item \textbf{相关课程}：编程与统计，生物信息学，遗传学，细胞生物学与遗传学，应用分子生物学
\end{itemize}

\section{实习经历}
\entry{上海宣泰医药科技有限公司}{2024.07 - 2024.09}{分析工程师 \textperiodcentered\ 分析部}{上海}
\begin{itemize}[parsep=0.2ex]
  \item \textbf{数据分析}：独立完成客户样本的前处理、HPLC 系统操作及数据分析，在严格的质控流程下确保检测结果的准确性和可靠性
  \item \textbf{质量控制}：通过严格的质控流程系统识别潜在杂质，为产品质量决策提供关键数据支持，展现了保证数据准确性与可追溯性的工程化思维
\end{itemize}

\section{科研与项目经历}

\entry{多模态大语言模型提取媒介性状数据评估：高精确率与结构依赖型记录召回}{2025.12 - 2026.08}{}{伦敦}
\role{Python、Gemma 4 31B IT（多模态LLM）、PyMuPDF、Docling、提示工程、JSON Schema 校验、二分图记录匹配评估、pytest}{}
\begin{itemize}[parsep=0.2ex]
  \item \textbf{基准构建}：设计加权多样性贪心算法从 VecTraits 超4万行、334篇论文中选取20篇论文，构建2,335条人工基准记录的开发集与431条留出测试集；参照 MIReVTD 标准将91字段的原始schema精简为22字段提取schema
  \item \textbf{多模态提取流程}：基于 PyMuPDF/Docling 构建文本与表格多模态证据包，驱动 Gemma 4 31B IT（256K 上下文多模态大语言模型）通过提示工程与 JSON Schema 约束完成可溯源性状记录提取；经11轮提示词迭代优化，开发集条件宏观字段 F1 达 0.810
  \item \textbf{效果评估与问题定位}：设计二分图记录匹配评估框架，按文档级记录遗漏与字段级差异对提取错误系统分类溯源；在留出测试集上评估锁定流程，取得记录精确率 1.000、召回率 0.360、条件宏观字段 F1 0.729，为大模型应用效果的持续优化提供问题定位依据
\end{itemize}

\entry{气候变化对佛罗里达州埃及伊蚊种群动态的影响建模}{2025.03 - 2025.06}{}{伦敦}
\role{R、deSolve、ggplot2、ODEs、统计验证（RMSE、Pearson、Spearman）}{}
\begin{itemize}[parsep=0.2ex]
  \item \textbf{模型适配与验证}：将已发表的热带 ODE 房室模型适配至亚热带地理环境，基于 53 周本地监测数据验证模型（RMSE=1.65，Pearson r=0.29，Spearman $\rho$=0.30），优化时间滞后与降雨尺度因子
\end{itemize}

\entry{生态建模与种群动力学模拟}{2024.01 - 2024.02}{}{伦敦}
\role{R、ggplot2、Lotka-Volterra、逻辑斯蒂增长模型}{}
\begin{itemize}[parsep=0.2ex]
  \item \textbf{建模项目组}：构建四个种群动力学模型——集合种群、渔业捕捞（MSY）、资源竞争及三物种捕食-被捕食模型，均使用差分方程与随机扩展实现；四项均获 90+ 评分（最高 94 分），展现扎实的算法实现与参数优化能力
\end{itemize}

\section{技能/其他}
\entrySkillsSep
\begin{itemize}[parsep=0.2ex]
  \item \textbf{AI/LLM 工程}：多模态大语言模型应用开发（Gemma 4 31B IT）、面向 NLP 信息抽取任务的提示工程与迭代优化、JSON Schema 结构化输出校验、二分图记录匹配评估与错误溯源分析、多后端 LLM 集成（Gemini API、Ollama 本地推理、OpenAI 兼容接口、HuggingFace 多模态模型）
  \item \textbf{Python 与开发工具}：Python（主力语言），Git，虚拟环境，pytest 测试，基于 JSON 配置/状态跟踪的异步子进程编排，Claude Code / Cursor 辅助开发
  \item \textbf{Web 开发}：FastAPI（路由/服务架构）、React、TypeScript、Vite、Tailwind CSS
  \item \textbf{数据分析与建模}：R（ggplot2、dplyr、caper、deSolve）、统计推断（GLM、PGLS、Wilcoxon 检验、t 检验、Pearson/Spearman 相关）、基于微分方程的动态系统建模
  \item \textbf{语言}：英语（雅思 7.5），普通话（母语），法语（A2）；\textbf{兴趣爱好}：钢琴，滑雪
\end{itemize}

\end{document}
